% Paper: a sealed ownership discipline in Lean 4.
%
% LIPIcs article, class lipics-v2021; build from this directory:
%   latexmk -pdf -interaction=nonstopmode sealed-ownership.tex
%
% Status note: this is a draft written against the implementation at the time of
% writing. Sections 3--7 describe checked Lean declarations; Section 8 (the C
% backend) and the parts marked *planned* are not implemented and are labelled as
% such in the text. Claims are tied to evidence in a table in Section 9.

\documentclass[a4paper,UKenglish,cleveref]{lipics-v2021}

\usepackage{booktabs}
\usepackage{tikz}
\usetikzlibrary{positioning,arrows,calc}

\setlength{\emergencystretch}{3em}
\tolerance=300

\makeatletter
\def\leansep{\penalty\z@}
\makeatother
\ExplSyntaxOn
\NewDocumentCommand{\lean}{m}{
  \tl_set:Nn \l_tmpa_tl {#1}
  \tl_replace_all:Nnn \l_tmpa_tl {\_} {\textunderscore\leansep}
  \tl_replace_all:Nnn \l_tmpa_tl {/} {/\leansep}
  \tl_replace_all:Nnn \l_tmpa_tl {.} {.\leansep}
  \texttt{\tl_use:N \l_tmpa_tl}
}
\ExplSyntaxOff

\DeclareUnicodeCharacter{2264}{\ensuremath{\leq}}
\DeclareUnicodeCharacter{2260}{\ensuremath{\neq}}
\DeclareUnicodeCharacter{2261}{\ensuremath{\equiv}}
\DeclareUnicodeCharacter{0393}{\ensuremath{\Gamma}}
\DeclareUnicodeCharacter{03B1}{\ensuremath{\alpha}}
\DeclareUnicodeCharacter{03B2}{\ensuremath{\beta}}
\DeclareUnicodeCharacter{03C1}{\ensuremath{\rho}}

\lstset{
  basicstyle=\ttfamily\small,
  columns=fullflexible,
  keepspaces=true,
  breaklines=true,
  breakatwhitespace=false,
  extendedchars=true,
  frame=leftline,
  framesep=3mm,
  rulecolor=\color{black!25},
  backgroundcolor=\color{black!1},
  xleftmargin=2mm,
  tabsize=2,
  literate=
    {→}{{$\rightarrow$}}2
    {∧}{{$\wedge$}}2
    {∀}{{$\forall$}}2
    {∃}{{$\exists$}}2
    {≤}{{$\leq$}}2
    {≠}{{$\neq$}}2
    {∈}{{$\in$}}2
    {⟨}{{$\langle$}}2
    {⟩}{{$\rangle$}}2
    {…}{{$\ldots$}}2
    {Γ}{{$\Gamma$}}2
    {α}{{$\alpha$}}2
    {β}{{$\beta$}}2
    {≡}{{$\equiv$}}2
    {¬}{{$\neg$}}2
}

\title{Borrowing Without Linear Types: A Sealed Ownership Discipline in Lean}
\titlerunning{A Sealed Ownership Discipline in Lean}
\author{Gavin Mendel-Gleason}{Scidonia Limited}{gavin@scidonia.ai}{}{}
\authorrunning{G. Mendel-Gleason}
\Copyright{Gavin Mendel-Gleason}
\DOIPrefix{}
\ccsdesc[500]{Theory of computation~Logic and verification}
\ccsdesc[300]{Software and its engineering~Software verification and validation}
\keywords{ownership, borrowing, Lean~4, permission calculi, verified systems}

\begin{document}
\maketitle

\begin{abstract}
Linear type systems forbid aliasing by construction. We ask whether the same
discipline can be enforced \emph{from outside} the host type system, by a checked
data structure, and connected to real mutation. We present a borrowing
discipline embedded in ordinary dependent type theory (Lean~4): handles are
freely copyable names with no authority of their own; a ledger of fractional
shares and byte footprints owns all authority; every operation is checked against
the current execution state. The discipline is realised by an executable checker
that is also the interpreter, with a machine-checked theorem that every accepted
execution preserves the world invariant, and by a second, result-returning
interface in which stale or refused operations are not even well-typed programs.
We report the current state of the implementation: a spatially and fractionally
splitting ledger with byte budgets, a command language with allocation, writes,
splits, joins, loans and disposal, sealed certificates with sequential
composition, a closed-run lifecycle, and union-find as a client, including a
path-compressing search that runs as a program in the discipline. The paper also
states precisely what is \emph{not} yet established: reads are observations, loans
are exclusive suspension rather than reborrowing, termination is fuel, module
encapsulation is not enforced by Lean, and the C backend with its representation
relation remains to be implemented. We argue that the proved core is enough to
make the sealed-ownership claim precise, and we separate the three trust levels
(abstract model, C implementation, compiled execution) that the full claim
requires.
\end{abstract}

\section{Introduction}

A verified functional specification with controlled low-level mutation is a
familiar tension. In Lean~4, arrays are persistent in their specification but
copy-on-write in their runtime: an unshared array is updated in place, a shared
one is copied. The observable behaviour is the same; the cost is not, and the
compiler's decision depends on reference counts that the type system does not
mention.

A complementary line of work studies that decision directly: static accounting
for reference-count-one updates, with proofs about the runtime representation.
This paper asks a different question. Suppose we do not modify the compiler at
all. Can we enforce an \emph{ownership discipline} in ordinary Lean — aliasing
rules, exclusive mutation, fractional reads, loans and disposal — such that a
program written against the discipline is checked, and its execution is licensed
by the checker? And can that be done without linear host types, in a language
where every value may be copied?

Our answer is a ledger. A \emph{handle} (allocation identifier, capability
identifier) is a name; names carry no authority. All authority lives in a
permission context that the operations consult. Copying a name is free and
grants nothing, because copying is not using: a duplicate name is just another
integer, and the next operation is checked against the current state. A stale
proof about an earlier state cannot authorise a live mutation.

\paragraph{Contributions.}
\begin{enumerate}
\item A permission calculus in Lean: shares over a context denominator,
  byte-footprint intervals, a decomposition tree with a byte-wise budget, a
  decidable write guard proved equal to its logical form, and a multi-owner
  context whose invariant (coverage, leaves within their allocation, distinct
  capability identifiers, one owner per allocation, storage agreement) is
  preserved by every transition (Sections~\ref{sec:calculus}--\ref{sec:world}).
\item An executable command language whose checker is its interpreter and a
  theorem that every accepted execution preserves the world invariant
  (Section~\ref{sec:unified}).
\item A result-returning, indexed program interface in which reads and
  data-dependent control flow compose; context-changing constructors carry their
  evidence, so refused operations are untypable rather than failing at run time
  (Section~\ref{sec:programs}).
\item Sealed certificates: programs packaged with their acceptance derivation,
  sequential composition through the actual intermediate world, and a closed-run
  lifecycle from the canonical empty world with snapshots as the only
  observation (Section~\ref{sec:sealed}).
\item Case studies: a spatial fill lifecycle, union-find through the checked
  route with a path-compressing search as a program, loans as suspension, and a
  byte-array backend refined against a pure model (Sections~\ref{sec:cases}).
\item An explicit map of the trust boundary: what is proved, what is
  assumed, and what the C backend would require (Sections~\ref{sec:c} and
  \ref{sec:limits}).
\end{enumerate}

The implementation is ordinary Lean with no axioms beyond the standard three
(\lean{propext}, \lean{Classical.choice}, \lean{Quot.sound}); an axiom audit
checks every headline theorem, and no proof uses \lean{sorry}. Statements quoted
below are the declarations in the development, with hypotheses shortened where
marked.

\section{Handles, authority, and the shape of the discipline}
\label{sec:design}

The design separates three things that linear types conflate:

\begin{itemize}
\item \textbf{identity}: allocation, capability, loan and lifetime identifiers
  are distinct types of tagged naturals (\lean{AllocId}, \lean{CapId},
  \lean{LoanId}, \lean{LifetimeId}); distinct types mean the checker cannot
  confuse an allocation for a capability;
\item \textbf{authority}: fractional shares over a context denominator, byte
  footprints, a ceiling and a capability state, stored in a ledger;
\item \textbf{state}: heap contents, kept separate from permissions, because
  contents change without permissions changing and vice versa.
\end{itemize}

A program is a sequence of operations over a \emph{world} (permission context,
backends, loans). The checker is the interpreter: an accepted execution returns
the world it executed. This is the sense in which the discipline is sealed at the
level of the pure model; Section~\ref{sec:limits} is explicit about what sealing
does and does not mean in Lean.

\section{The permission calculus}
\label{sec:calculus}

\paragraph{Shares.} A capability's share of a resource is a fraction $q/D$: a
numerator over a context-wide denominator. Only the whole share $q=D$ may write;
splits produce shares that add to exactly the parent's share; hence two partial
shares can never both write. The arithmetic is not a convention: the write guard
requires \lean{c.q = D}, and the budget invariant below makes $q = D$
achievable by at most one capability at a time.

\paragraph{Footprints.} A footprint is a half-open byte interval $[lo,hi)$ of
one allocation. Splits cut an interval into two; disjointness, adjacency and
coverage are proved.

\paragraph{The ledger.} One \lean{Owner} per allocation holds a tree of leaves.
A leaf is a capability: footprint, share numerator, ceiling, state
(\lean{active}, \lean{retired}, \lean{consumed}). Internal nodes are
\emph{spatial} (children split the byte range) or \emph{fractional} (children
split the share over the same range). The key measure is the \emph{budget}: at a
byte $x$, the sum of shares of all leaves covering $x$. Coverage is:

\begin{lstlisting}
def Owner.Covered (o : Owner) (D : Nat) : Prop :=
  ∀ x, x < o.len → o.root.budget x = D
\end{lstlisting}

Every byte is exactly accounted for. Splits preserve the budget by
construction; the sibling-only fractional join merges children of summed share
and preserves it, and is refused when the children carry the same capability
identifier.

The write guard is decidable, and provably the same predicate the proofs use:

\begin{lstlisting}
def Owner.writableB (o : Owner) (D : Nat) (id : CapId) (i : Nat) : Bool :=
  match o.root.find? id with
  | none => false
  | some c =>
    decide (c.state = .active ∧ c.q = D ∧ c.ceiling = .readWrite
              ∧ c.lo ≤ i ∧ i < c.hi)

theorem Owner.writableB_sound :
    o.writableB D id i = true → o.Writable D id i
\end{lstlisting}

That small theorem is load-bearing: the executable check and the logical guard
cannot drift.

\section{Contexts and the world}
\label{sec:world}

A \lean{Ctx} is a list of owners, the denominator, and a fresh-allocation
counter. \lean{Ctx.Realizes} requires every owner to be covered and allocation
identifiers to be distinct. The stronger \lean{Ctx.WellFormed} adds that every
owner's leaves lie inside its own allocation, and that capability identifiers
across all owners are distinct.

Freshness is enforced, not assumed: both split dispatch and the join require the
new identifiers to be absent from the whole context, and allocation requires its
root and disposal identifiers fresh and distinct. This is what stops a new
capability from colliding with an old one and inheriting its authority.

The world adds the byte backends and the loan bookkeeping, with the bundle

\begin{lstlisting}
def World.WellFormed (w : World) : Prop :=
  w.ctx.WellFormed           -- coverage, leaves within, identifiers distinct
  ∧ w.ctx.Functional         -- one owner per allocation
  ∧ w.StorageMatches         -- the found owner's backend has its length
\end{lstlisting}

\lean{World.update\_storage} is the workhorse: replacing one owner by a
same-length result keeps the heap agreement, so every command's preservation
proof ends there.

\section{The unified checker and interpreter}
\label{sec:unified}

The executable language is a command list; the checker folds it over the world:

\begin{lstlisting}
inductive UCmd where
  | write (a : AllocId) (id : CapId) (i : Nat) (v : UInt8)
  | splitAt (a : AllocId) (id id1 id2 : CapId) (k : Nat)
  | splitFraction (a : AllocId) (id id1 id2 : CapId) (q1 q2 : Nat)
  | join (a : AllocId) (node newId : CapId)
  | free (a : AllocId) (d : CapId)
  | alloc (n : Nat) (rootId disposal : CapId)
  | beginLoan (a : AllocId) (lender : CapId) (loanId : LoanId)
  | endLoan (a : AllocId) (loanId : LoanId)
  | nop
\end{lstlisting}

One theorem covers every command:

\begin{lstlisting}
theorem checkWorld_wf :
    checkWorld w cs = .ok w' → w.WellFormed → w'.WellFormed
\end{lstlisting}

Accepted execution preserves coverage, freshness, owner uniqueness and storage
agreement — through allocation, writes, spatial and fractional splits, joins,
loans and disposal. The older relation \lean{World.Rel} follows from it.

\paragraph{Loans are suspension.} Beginning a loan retires the lender's leaf.
Because the write guard requires an \lean{active} leaf, a write through the
lender is then refused by the \emph{same} guard — there is no parallel loan
machinery to keep in sync. Ending the loan restores the leaf. The refusal is a
theorem with a positive control: the write is accepted before the loan, refused
while it is outstanding, and accepted again after restoration.

\section{Programs with post-contexts}
\label{sec:programs}

The command language is a trace language: a fixed list of operations. An
algorithm with data-dependent control flow — union-find's search, which reads a
parent and decides where to go next — cannot be a list. We therefore add two
interfaces.

\paragraph{The context-preserving fragment.} \lean{Owned Γ α} is reads, writes
and monadic binding over a fixed permission context. The fuelled evaluator is
total, with exhaustion a visible \lean{none}. Two theorems: a successful run
preserves the world invariant, and it leaves the context unchanged, realising the
fragment's constant index.

\paragraph{The indexed interface.} \lean{IProg Γ α Q} starts with context
\lean{Γ}, returns \lean{a : α} and ends with \lean{Q a}. Context-changing
constructors carry their evidence: allocation needs fresh identifiers, a spatial
split needs the ledger transition to succeed, disposal needs the disposal check.
A refused operation is therefore not a program at all. For example, the
freshness evidence for reusing a parent capability as a child is uninhabited:

\begin{lstlisting}
theorem stale_child_refused :
    ¬ (Owner.init ⟨1⟩ 4 64 (by decide) ⟨3⟩ ⟨4⟩).FreshIds ⟨3⟩ ⟨5⟩
\end{lstlisting}

The evaluator is total by fuel, and the layer's soundness theorem is:

\begin{lstlisting}
theorem evalIProg_post :
    w.ctx = Γ → evalIProg fuel p w = some (a, w') → w.WellFormed →
    w'.ctx = Q a ∧ w'.WellFormed
\end{lstlisting}

Execution from a world matching the typing context ends in a world matching the
postcondition, with the invariant preserved. Post-contexts are computed with
\lean{Option.getD} so the type function stays computable while the evidence
records success.

\paragraph{Recursion and loops.} Recursion does not create a new loan: a
recursive call can sequentially use the existing exclusive capability while its
caller is suspended on the stack. Because the *host* recursion builds the program
and the evaluator runs it, a fuel-bounded search is an ordinary program:

\begin{lstlisting}
def findCompress (Γ : Ctx) (a : AllocId) (id : CapId) :
    Nat → Nat → Owned Γ (Except Unit Nat)
  | 0, _ => .pure (.error ())
  | fuel + 1, x =>
    .bind (.read a id x) fun
      | none => .pure (.error ())
      | some p =>
        if p.toNat = x then .pure (.ok x)
        else .bind (findCompress Γ a id fuel p.toNat) fun
          | .error e => .pure (.error e)
          | .ok root => .bind (.write a id x root.toUInt8)
                          fun _ => .pure (.ok root)
\end{lstlisting}

On the four-node union-find world with parents \lean{[0,0,0,2]}, the search from
node 3 returns root 0 and compresses the path to \lean{[0,0,0,0]}, with the final
world well formed and its context unchanged. Safety (every access permitted,
invariant preserved) and functional correctness (the root returned, the
compression written) are separate theorems; termination is the fuel, and for
valid forests a termination theorem is future work.

\section{Sealed execution}
\label{sec:sealed}

\paragraph{Certificates.} \lean{Prog pre} packages a command list with the proof
that the checker accepts it from \lean{pre}; it cannot exist without the
derivation. Sequential composition runs an appended list from the first part's
result, so the second program provably starts in the first's post-world.

\paragraph{Closed runs.} \lean{ClosedProgram} packages commands accepted from the
canonical empty world, whose only observation is a snapshot of the final world.
The indexed interface has the matching closed form: a program that allocates,
writes and disposes, typed as ending with no owners, accepted, well formed, with
the disposed allocation's backend gone.

\paragraph{Failure.} Allocation failure, explicit algorithm errors and
nontermination are modelled separately: the evaluator's fuel exhaustion and the
checker's errors are visible outcomes, not silent defaults. Resource-preserving
failure results for the C backend are part of the planned work
(Section~\ref{sec:c}).

\section{Case studies}
\label{sec:cases}

\paragraph{A spatial and fractional lifecycle.} A program allocates a fresh
4-byte owner, writes through the root capability, splits spatially, writes the
left half, splits that half fractionally into two half-shares, joins them back
into a whole share, writes again, and disposes. The fractional split leaves two
half-shares, neither writable; the join restores the whole share and the second
write is accepted again. The lifecycle is accepted, ends well formed, and leaves
a pre-existing 64-byte backend untouched.

\paragraph{Union-find.} We treat union-find as the \emph{witness} for the
discipline rather than as a correctness target: it is the algorithm whose cost
distinguishes unique from shared access. Three layers exist. (i) The functional
model: reachability, \lean{Same} as an equivalence, simple union, the concrete
partition \lean{[0,0,0,2]} with every pair \lean{Same}, bounded root finding and
one-step compression with preservation of \lean{Same}. (ii) The checked trace:
unions issued through \lean{World.writeChecked} land on a backend whose parent
array is exactly the functional model's \lean{[0,0,0,2]}, and the world remains
well formed. (iii) The program: the path-compressing search of
Section~\ref{sec:programs} runs in the context-preserving fragment.
\emph{Planned}: the representation relation between world bytes and the
functional parent array, the forest invariant (parent bounds, valid forest) with
its preservation, and the general theorem that the program's result agrees with
the functional finder. Today these are separate; the representation layer is the
next step.

\paragraph{Backends.} Two implementations share one model. The pure backend is a
byte list with a length proof; the byte-array backend is the same interface over
Lean's \lean{ByteArray}. The refinement (empty, reads, writes, range fills) is
proved, and the 64-byte two-range fill through the byte-array route produces the
reference bytes. The tempting sentence — ``the byte array is updated in place''
— is exactly what is \emph{not} observable from inside, and we do not claim it.

\section{The C backend: boundary and trust}
\label{sec:c}

The sealed-ownership claim ultimately concerns real mutation. The planned
architecture retains Lean control flow and supplies a narrow C interface for
allocation, validated reads and writes, snapshots and disposal, with no raw
pointers as ordinary client values:

\begin{lstlisting}[language=C]
/* Design sketch, not a Lean FFI declaration. */
status buffer_alloc(arena *, size_t count, buffer_id *out);
status buffer_read(arena *, buffer_id, size_t index, word *out);
status buffer_write(arena *, buffer_id, size_t index, word value);
status buffer_free(arena *, buffer_id);
\end{lstlisting}

The proof obligation is a representation relation \lean{BackendRel concrete
model} and, per primitive, that every permitted model execution is implemented by
the concrete step and the relation is preserved: extent, layout, alignment,
index arithmetic and overflow; live identity, generation and pointer validity;
initialised reads and permitted writes; no use after disposal and no double
disposal; confinement of inputs and outputs and independent immutable snapshots;
allocation failure and error cleanup; and the sequential-execution assumption,
including callbacks and reentrancy.

Three levels of evidence must stay distinct, and only the first exists today:
(1) the abstract backend model proves the intended primitives preserve ownership
and observations; (2) the implemented C primitives satisfy those contracts under
a chosen semantics or explicit assumptions — \emph{not implemented}; (3) the
compiler, runtime and ABI preserve behaviour — either proved or named as trusted.
A Lean theorem about an abstract external operation does not verify the C body.
The manuscript must say whether the C implementation is verified,
translation-validated, or trusted; tests support engineering confidence but do
not replace that distinction. Two implementation routes are available: Lean
control flow with C primitives (recommended first), or a custom lowering of the
fragment with a simulation proof (additional compiler work).

\section{Evaluation and evidence ledger}
\label{sec:eval}

What is mechanically checked today (all by the repository's gate: builds, axiom
audit admitting only the standard three, and the paper/document builds):

\begin{center}\small
\begin{tabular}{@{}p{0.46\linewidth}p{0.44\linewidth}@{}}
\toprule
Claim & Evidence \\
\midrule
Checker and interpreter are one; accepted runs preserve the invariant &
\lean{checkWorld\_wf}, \lean{checkWorld\_rel} \\
Refused operations are untypable in the indexed interface &
\lean{stale\_child\_refused}, evidence-carrying constructors \\
Execution realises the typed post-context &
\lean{evalIProg\_post} \\
Closed runs start from nothing and end with no owners &
\lean{iClosed\_ok}, \lean{iClosed\_closed}, \lean{iClosed\_freed} \\
Union-find's compression search is a program; result and compression proved &
\lean{findCompress\_run\_root}, \lean{findCompress\_compresses} \\
The checked union trace lands on the functional parent array &
\lean{ufTrace\_backend}, \lean{ufTrace\_partition} \\
Byte-array backend refines the pure model &
\lean{empty\_refines}, \lean{write\_refines}, \lean{ba64\_snapshot} \\
Loans: suspended write refused, restored write accepted &
\lean{uLoan\_suspended\_write\_refused}, \lean{uLoanRun\_ok} \\
\midrule
Representation relation world$\leftrightarrow$array, forest invariant & planned \\
General find/union correctness with compression & planned \\
C primitives and their refinement & planned \\
Compiled performance evidence & planned \\
\bottomrule
\end{tabular}
\end{center}

\paragraph{Timing.} One interpreted observation on the 64-byte fill-and-sum
pattern, with agreeing checksums: 73,959~ns per repetition for the pure list
backend, 47,688 for the byte-array backend, 51,956 for the reference heap. Because
the run is interpreted, these numbers include the evaluator's own overhead; they
are environment-dependent and say nothing about compiled execution, in-place
update, or a real runtime. We report them as an observation, not a claim.

\section{Limitations and threats to validity}
\label{sec:limits}

\begin{itemize}
\item \textbf{Reads are observations.} The discipline constrains writes. A read
  through an unheld capability is not refused, and fractional read permissions
  are not yet distinguished from free observation. This is the largest gap on the
  ownership side.
\item \textbf{Loans are exclusive suspension.} The borrower is the lender's own
  identifier; there is no reborrow capability, no non-root lifetime in the
  executable language, and the policy-level liveness conjunction is not consulted
  by the checker.
\item \textbf{Termination is fuel.} Programs are total by fuel; exhaustion is a
  visible outcome. Termination for valid forests and for the general loops is
  future work.
\item \textbf{No module opacity.} Lean has no module boundary that hides a
  definition; the seal is by discipline, not enforced by the language. A safe
  result type and an opaque runtime are future work.
\item \textbf{One invariant, not the whole design invariant.} Coverage, leaves
  within, distinct identifiers, owner uniqueness and storage agreement are
  maintained. Provenance roles of internal split nodes, identifier recycling
  policy after disposal, and loan-to-capability relationships are not yet part of
  a preserved invariant.
\item \textbf{Union-find correctness is partial.} The representation layer and
  the general partition correspondence are planned; today's theorems concern the
  concrete four-node instance and the search program.
\item \textbf{No verification of C.} The backend exists as a Lean model and a
  ByteArray refinement only.
\end{itemize}

\section{Related work}

Borrow checking in Rust enforces aliasing rules in the compiler; we enforce an
analogous discipline in a proof assistant without changing the compiler.
Linear and uniqueness type systems (Linear Haskell, Clean's uniqueness types)
make the discipline part of typing; our handles are freely copyable and the
discipline lives in state. Separation logic and its implementations (Iris,
RefinedC) reason about heap ownership with resource algebras very much like our
shares and footprints; we differ in target (an executable, sealed runner inside
Lean) and in the explicit trust boundary with C. Verified compilers (CompCert)
prove preservation of whole-program semantics; our claim is an ownership
discipline plus a boundary specification, and we do not claim compiler
verification. The companion work in this repository studies static
reference-count-one accounting for in-place updates at the compiler level; this
paper is the dual: the discipline that a runtime checker can enforce when the
compiler is unchanged.

\section{Conclusion}

A borrowing discipline need not be a type system. Handles as names, a ledger of
fractional shares and footprints, and a checker that is also the interpreter give
an ownership discipline in ordinary Lean, with machine-checked preservation and
a program interface whose refused operations do not typecheck. The remaining work
is well-defined: a representation relation and forest invariant for union-find,
fractional reads and real reborrowing, module-level sealing, and the C backend
with its three trust levels made explicit.

\begin{thebibliography}{9}
\bibitem{lean4} L. de Moura and S. Ullrich. The Lean~4 theorem prover and
programming language. \emph{Automated Deduction -- CADE 28}, 2021.
\bibitem{rust} S. Klabnik and C. Nichols. \emph{The Rust Programming Language}.
No Starch Press, 2019.
\bibitem{linearhaskell} J.-P. Bernardy, M. Boespflug, R. Newton, S. Peyton Jones,
and A. Spiwack. Linear Haskell: practical linearity in a higher-order polymorphic
language. \emph{POPL 2018}.
\bibitem{clean} E. Barendsen and S. Smetsers. Uniqueness typing for functional
languages with graph rewriting. \emph{Term Graph Rewriting}, 1993.
\bibitem{iris} R. Jung, R. Krebbers, J.-H. Jourdan, A. Bizjak, L. Birkedal, and
D. Dreyer. Iris from the ground up. \emph{JFP} 28, 2018.
\bibitem{refinedc} M. Sammler et al. RefinedC: automating the foundational
verification of C code with refined ownership types. \emph{PLDI 2021}.
\bibitem{compcert} X. Leroy. Formal verification of a realistic compiler.
\emph{CACM} 52(7), 2009.
\bibitem{companion} G. Mendel-Gleason. Ownership for in-place unobservable
updates in Lean~4. Companion manuscript, this repository.
\end{thebibliography}

\end{document}
